\begin{center}
    {\LARGE \textbf{2015物理化学}} \\[2ex]
    {\large 时量180分钟 \quad 满分150分}
\end{center}

\vspace{0.5cm}
\noindent\textbf{注意事项:}
\begin{enumerate}[label=\arabic*., parsep=0pt, itemsep=0pt]
    \item 答题前,考生先将自己的姓名、学号、考场号、座位号填写在答题卡上,将条形码准确粘贴在考生信息条形码粘贴区。
    \item 考试结束后,将本试卷与答题纸一并收回。
    \item 回答各个问题时,将答案写在答题纸上,写在本试卷上无效。
\end{enumerate}

\vspace{0.5cm}
\noindent\textbf{一、选择题（30 pts.）}

\begin{enumerate}[label=\arabic*., itemsep=1.5ex]
    \item 一定量的理想气体从同一始态出发，分别经（1）等温压缩，（2）绝热压缩到具有相同压力的终态，以$H_1$、$H_2$分别表示两个终态的焓值，则有
    \begin{tasks}(4)
        \task $H_1 > H_2$
        \task $H_1 = H_2$
        \task $H_1 < H_2$
        \task 不确定
    \end{tasks}
    \item 用$130^\circ\mathrm{C}$的水蒸气（蒸气压为$2.7p^\ominus$）与$1200^\circ\mathrm{C}$的焦炭反应生成水煤气：$\ce{C + H2O -> CO + H2}$，如果通入的水蒸气反应掉$70\%$，问反应后混合气体中$\ce{CO}$的分压为多少？（设总压$2.7p^\ominus$不变）
    \begin{tasks}(4)
        \task $0.72p^\ominus$
        \task $1.11p^\ominus$
        \task $2.31p^\ominus$
        \task $1.72p^\ominus$
    \end{tasks}
    \item $\ce{CaCO3(s)}$、$\ce{CaO(s)}$、$\ce{BaCO3(s)}$、$\ce{BaO(s)}$及$\ce{CO2(g)}$构成的一个平衡物系，其组分数为
    \begin{tasks}(4)
        \task $2$
        \task $3$
        \task $4$
        \task $5$
    \end{tasks}
    \item 某电池反应为$\ce{Zn(s) + Mg^{2+}(a=0.1) = Zn^{2+}(a=1) + Mg(s)}$，用实验测得该电池的电动势$E=0.231\ \mathrm{V}$，则电池的$E^\ominus$为
    \begin{tasks}(4)
        \task $0.2903\ \mathrm{V}$
        \task $-0.2312\ \mathrm{V}$
        \task $0.0231\ \mathrm{V}$
        \task $-0.2022\ \mathrm{V}$
    \end{tasks}
    \item 理想气体在等温条件下反抗恒定外压膨胀，该变化过程中体系的熵变$\Delta S_\text{体}$及环境的熵变$\Delta S_\text{环}$应为
    \begin{tasks}(2)
        \task $\Delta S_\text{体}>0,\ \Delta S_\text{环}=0$
        \task $\Delta S_\text{体}<0,\ \Delta S_\text{环}=0$
        \task $\Delta S_\text{体}>0,\ \Delta S_\text{环}<0$
        \task $\Delta S_\text{体}<0,\ \Delta S_\text{环}>0$
    \end{tasks}
    \item 在$373.15\ \mathrm{K}$时，某有机液体$\ce{A}$和$\ce{B}$的蒸气压分别为$p$和$3p$，$\ce{A}$和$\ce{B}$的某混合物为理想液体混合物，并在$373.15\ \mathrm{K}$、$2p$时沸腾，那么$\ce{A}$在平衡蒸气相中的摩尔分数是多少？
    \begin{tasks}(4)
        \task $1/3$
        \task $1/4$
        \task $1/2$
        \task $3/4$
    \end{tasks}
    \item 已知$298\ \mathrm{K}$时，下列电极电势：
    $\varphi^\ominus(\ce{Zn^{2+}, Zn})=-0.7628\ \mathrm{V}$，$\varphi^\ominus(\ce{Cd^{2+}, Cd})=-0.4029\ \mathrm{V}$，$\varphi^\ominus(\ce{I2, I^-})=0.5355\ \mathrm{V}$，$\varphi^\ominus(\ce{Ag^+, Ag})=0.7991\ \mathrm{V}$，下列电池的标准电动势最大的是
    \begin{tasks}(1)
        \task $\ce{Zn(s) | Zn^{2+} \parallel Cd^{2+} | Cd(s)}$
        \task $\ce{Zn(s) | Zn^{2+} \parallel H^+ | H2, Pt}$
        \task $\ce{Zn(s) | Zn^{2+} \parallel I^- | I2, Pt}$
        \task $\ce{Zn(s) | Zn^{2+} \parallel Ag^+ | Ag(s)}$
    \end{tasks}
    \item 有三种电极表示式：
    (1) $\ce{Pt, H2(p^\ominus) | H+(a=1)}$，(2) $\ce{Cu | Pt, H2(p^\ominus) | H+(a=1)}$，(3) $\ce{Hg | Cl^-(aq) | H2(p^\ominus) | H+(a=1)}$，则氢电极的电极电势彼此关系为
    \begin{tasks}(4)
        \task 逐渐变大
        \task 逐渐变小
        \task 不能确定
        \task 彼此相等
    \end{tasks}
    \item 恒温恒压条件下，某化学反应若在电池中可逆进行时吸热，据此可以判断下列热力学量中何者一定大于零？
    \begin{tasks}(4)
        \task $\Delta U$
        \task $\Delta H$
        \task $\Delta S$
        \task $\Delta G$
    \end{tasks}
    \item 组分$\ce{A}$和$\ce{B}$可以形成四种稳定化合物：$\ce{A2B}$，$\ce{AB}$，$\ce{AB2}$，$\ce{AB3}$，设所有这些化合物都有相合熔点，则此体系的低共熔点最多有几个？
    \begin{tasks}(4)
        \task $3$
        \task $4$
        \task $5$
        \task $6$
    \end{tasks}
    \item 在一般情况下，电位梯度只影响
    \begin{tasks}(2)
        \task 离子的电迁移率
        \task 离子迁移速率
        \task 电导率
        \task 离子的电流分数
    \end{tasks}
    \item 当一反应物的初始浓度为$0.04\ \mathrm{mol\cdot dm^{-3}}$时，反应的半衰期为$360\ \mathrm{s}$，初始浓度为$0.024\ \mathrm{mol\cdot dm^{-3}}$时，半衰期为$600\ \mathrm{s}$，此反应为
    \begin{tasks}(4)
        \task $0$级反应
        \task $1.5$级反应
        \task $2$级反应
        \task $1$级反应
    \end{tasks}
    \item 对于$1$-$1$型电解质溶液的摩尔电导率可以看作是正、负离子的摩尔电导率之和，这一规律只适用于
    \begin{tasks}(2)
        \task 强电解质
        \task 弱电解质
        \task 无限稀释电解质溶液
        \task 摩尔浓度为$1$的溶液
    \end{tasks}
    \item 反应$\ce{2N2O5 -> 4NO2 + O2}$的速率常数单位是$\mathrm{s^{-1}}$。对该反应的下述判断哪个对？
    \begin{tasks}(4)
        \task 单分子反应
        \task 双分子反应
        \task 复合反应
        \task 不能确定
    \end{tasks}
    \item 对于气相基元反应，按过渡态理论，不正确的关系式是
    \begin{tasks}(2)
        \task $E_\mathrm{a} = \Delta^\neq U_\mathrm{m}^\ominus + RT$
        \task $E_\mathrm{a} = \Delta^\neq H_\mathrm{m}^\ominus + nRT$
        \task $E_\mathrm{a} = E_0 + RT$
        \task $E_\mathrm{a} = E_0 + mRT$
    \end{tasks}
\end{enumerate}

\vspace{0.5cm}
\noindent\textbf{二、计算题（20 pts.）}

在$101.325\ \mathrm{kPa}$下，甲苯的沸点为$383.2\ \mathrm{K}$，已知此条件下$\Delta S_\mathrm{m}(\mathrm{l})=2619.1\ \mathrm{J\cdot K^{-1}\cdot mol^{-1}}$，$\Delta_\mathrm{vap}S_\mathrm{m}=87.03\ \mathrm{J\cdot K^{-1}\cdot mol^{-1}}$。试计算：
\begin{enumerate}[label=(\arabic*), itemsep=1ex]
    \item 上述条件下$1\ \mathrm{mol}$甲苯(l)蒸发时的$Q$，$W$，$\Delta_\mathrm{vap}U_\mathrm{m}$，$\Delta_\mathrm{vap}H_\mathrm{m}$，$\Delta_\mathrm{vap}A_\mathrm{m}$，$\Delta_\mathrm{vap}G_\mathrm{m}$及$S_\mathrm{m}(\mathrm{g})$。
    \item 欲使甲苯在$373.2\ \mathrm{K}$时沸腾，其相应的饱和蒸气压为多少？
\end{enumerate}

\vspace{0.5cm}
\noindent\textbf{三、计算题（20 pts.）}

$\ce{Mg}$(熔点$924\mathrm{K}$)和$\ce{Zn}$(熔点$692\mathrm{K}$)的相图具有两个低共熔点，一个为$641\ \mathrm{K}$($\ce{Mg}$的摩尔分数为$0.018$)，另一个为$620\ \mathrm{K}$($\ce{Mg}$的摩尔分数为$0.7212$)，在体系的熔点曲线上有一个最高点$863\ \mathrm{K}$($\ce{Mg}$的摩尔分数为$0.3339$)。
\begin{enumerate}[label=(\arabic*), itemsep=1ex]
    \item 绘出$\ce{Mg}$和$\ce{Zn}$的$T-x$图，并标明各区中的相。
    \item 分别指出含$\ce{Mg}$的摩尔分数为$0.5357$和$0.9150$的两个混合物从$973\ \mathrm{K}$冷却到$573\ \mathrm{K}$步冷过程中的相变，并根据相律予以说明。
    \item 绘出含$\ce{Mg}$的摩尔分数为$0.7212$的熔化物的步冷曲线。
\end{enumerate}

\vspace{0.5cm}
\noindent\textbf{四、计算题（15 pts.）}

$\ce{CO2}$在高温下按下式分解$\ce{2CO2 -> 2CO + O2}$，在$p^\ominus$、$1000\ \mathrm{K}$时解离度为$2.0\times10^{-7}$，$1400\ \mathrm{K}$时为$1.27\times10^{-4}$，设在该温度范围内反应热效应不随温度而改变，则在$1000\ \mathrm{K}$时反应的$\Delta_\mathrm{r}G_\mathrm{m}^\ominus$和$\Delta_\mathrm{r}S_\mathrm{m}^\ominus$各为多少？

\vspace{0.5cm}
\noindent\textbf{五、计算题（15 pts.）}

计算$1200^\circ\mathrm{C}$时，固态$\ce{FeO}$上氧的平衡压力，反应$\ce{FeO(s) <=> Fe(s) + \frac{1}{2}O2(g)}$在$25^\circ\mathrm{C}$时的$\Delta_\mathrm{r}H_\mathrm{m}^\ominus=268.82\ \mathrm{kJ\cdot mol^{-1}}$，且已知：

\begin{center}
\begin{tabular}{lccc}
 & $\ce{FeO}$ & $\ce{Fe}$ & $\ce{O2}$ \\
\hline
$S_\mathrm{m}^\ominus(298\ \mathrm{K})/(\mathrm{J\cdot K^{-1}\cdot mol^{-1}})$ & $58.79$ & $27.15$ & $205.03$ \\
$C_{p,\mathrm{m}}/(\mathrm{J\cdot K^{-1}\cdot mol^{-1}})$ & $57.70$ & $37.66$ & $33.68$ \\
\end{tabular}
\end{center}

\vspace{0.5cm}
\noindent\textbf{六、计算题（15 pts.）}

在$671\sim768\ \mathrm{K}$之间，$\ce{C2H5Cl}$气相分解反应（$\ce{C2H5Cl -> C2H4 + HCl}$）为一级反应，速率常数$k(\mathrm{s^{-1}})$和温度$T$的关系式为
\[
\lg(k/\mathrm{s^{-1}}) = -13290/(T/\mathrm{K}) + 14.6
\]
\begin{enumerate}[label=(\arabic*), itemsep=1ex]
    \item 求$E_\mathrm{a}$和$A$；
    \item 在$700\ \mathrm{K}$时，将$\ce{C2H5Cl}$通入一个反应器中（$\ce{C2H5Cl}$的起始压力为$26664.5\ \mathrm{Pa}$），反应开始后，反应器中压力增大，问需多少时间反应器中压力变为$46662.8\ \mathrm{Pa}$？
\end{enumerate}

\vspace{0.5cm}
\noindent\textbf{七、计算题（15 pts.）}

在$273\sim318\ \mathrm{K}$范围内，电池的电动势与温度的关系可由下列公式表示：

(1) $\ce{Cu(s) | Cu2O(s) | NaOH(aq) | HgO(s) | Hg(l)}$
\[
E = \{461.7 - 0.144(T/K - 298) + 0.00014(T/K - 298)^2\}\ \mathrm{mV}
\]

(2) $\ce{Pt, H2(p^\ominus) | NaOH(aq) | HgO(s) | Hg(l)}$
\[
E = \{925.65 - 0.2948(T/K - 298) + 0.00049(T/K - 298)^2\}\ \mathrm{mV}
\]
已知$\Delta_\mathrm{f}H_\mathrm{m}^\ominus(\ce{H2O,l})=-285.85\ \mathrm{kJ\cdot mol^{-1}}$，$\Delta_\mathrm{f}G_\mathrm{m}^\ominus(\ce{H2O,l})=-237.19\ \mathrm{kJ\cdot mol^{-1}}$，试分别计算$\ce{HgO(s)}$和$\ce{Cu2O(s)}$在$298\ \mathrm{K}$时的$\Delta_\mathrm{f}G_\mathrm{m}^\ominus$和$\Delta_\mathrm{f}H_\mathrm{m}^\ominus$。

\vspace{0.5cm}
\noindent\textbf{八、计算题（20 pts.）}

电池$\ce{Cd(s) | Cd(OH)2(s) | NaOH(0.01\ mol\cdot kg^{-1}) | H2(p^\ominus), Pt}$在$298\ \mathrm{K}$时电动势$E=0.000\ \mathrm{V}$，$(\partial E/\partial T)_p=0.002\ \mathrm{V\cdot K^{-1}}$，$\varphi^\ominus(\ce{Cd^{2+} | Cd})=-0.403\ \mathrm{V}$。
\begin{enumerate}[label=(\arabic*), itemsep=1ex]
    \item 写出电极反应和电池反应；
    \item 求电池反应的$\Delta_\mathrm{r}G_\mathrm{m}$，$\Delta_\mathrm{r}H_\mathrm{m}$，$\Delta_\mathrm{r}S_\mathrm{m}$；
    \item 求$\ce{Cd(OH)2}$的溶度积常数$K_{\mathrm{sp}}$。
\end{enumerate}